\documentclass[11pt,a4paper]{article}
			\usepackage[french]{babel}
					
				\usepackage{pifont}  
				\usepackage[utf8x]{inputenc}
				\usepackage[T1]{fontenc} 
				\usepackage{lmodern}			
				\usepackage{fancyhdr}
				\usepackage{textcomp}
				\usepackage{makeidx}
				\usepackage{tabularx}
				\usepackage{multicol}
				\usepackage{multirow}
				\usepackage{longtable}
				\usepackage{color}
				\usepackage{soul}
				\usepackage{boxedminipage}
				\usepackage{shadow}
				\usepackage{framed}			
				\usepackage{array}
				\usepackage{url}
				\usepackage{ragged2e}
				\usepackage{fancybox}
				\newcommand{\cadretitre}[2]{
				  \vspace*{0.8\baselineskip}
				  \begin{center}%
				  \boxput*(0,1){%
					%\colorbox{white}{\Large\textbf{\ #1\ }}%
				  }%
				  {%
					\setlength{\fboxsep}{10pt}%
				    \Ovalbox{\begin{minipage}{.8\linewidth}\begin{center}\Large\sffamily{#2}\end{center}\end{minipage}}}%
				  \end{center}
				  \vspace*{2\baselineskip}
				  }
			
			\makeatletter
			\def\@seccntformat#1{\protect\makebox[0pt][r]{\csname the#1\endcsname\quad}}
			\makeatother

				% Permet d'afficher qqchose à une positin absolue
				\usepackage[absolute]{textpos}
				\setlength{\TPHorizModule}{1cm}
				\setlength{\TPVertModule}{\TPHorizModule}
	
				\usepackage[titles]{tocloft}
				\setlength{\cftbeforesecskip}{0.5ex}
				\setlength{\cftbeforesubsecskip}{0.2ex}
				\addto\captionsfrench{\renewcommand\contentsname{}}
				
				\usepackage[font=scriptsize]{caption}
				
				\usepackage{listings}
\lstdefinestyle{lstverb}
  {
    basicstyle=\footnotesize,
    frameround=tttt, frame=trbl, framerule=0pt, rulecolor=\color{gray},
    lineskip=-1pt,   % pour rapprocher les lignes
    flexiblecolumns, escapechar=\\,
    tabsize=4, extendedchars=true
  }
\lstnewenvironment{Java}[1][]{\lstset{style=lstverb,language=java,#1}}{}
				\ifx\pdfoutput\undefined
					\usepackage{graphicx}
				\else
					\usepackage[pdftex]{graphicx}
				\fi
				\usepackage[a4paper, hyperfigures=true, colorlinks, linkcolor=black, citecolor=blue,urlcolor=blue, pagebackref=true, bookmarks=true, bookmarksopen=true,bookmarksnumbered=true,
                pdfauthor={}, pdftitle={TD Séquentiel - Rappels de base}, pdfkeywords={TD S\'equentiel - Rappels de base, },pdfpagemode=UseOutlines,pdfpagetransition=Dissolve,nesting=true,
				backref, pdffitwindow=true, bookmarksnumbered=true]{hyperref}
				\usepackage{supertabular}
				\usepackage[table]{xcolor}
				\usepackage{url}
				\usepackage{caption} 
				\setlength{\parskip}{1.3ex plus 0.2ex minus 0.2ex}
				\setlength{\parindent}{0pt}
				
				\makeatletter
				\def\url@leostyle{ \@ifundefined{selectfont}{\def\UrlFont{\sf}}{\def\UrlFont{\footnotesize\ttfamily}}}
				\makeatother
				\urlstyle{leo}
				
				\definecolor{examplecolor}{rgb}{0.156,0.333,0.443}
				\definecolor{definitioncolor}{rgb}{0.709,0.784,0.454}
				\definecolor{exercisecolor}{rgb}{0.49,0.639,0}
				\definecolor{hintcolor}{rgb}{0.941,0.674,0.196}
				\definecolor{tableHeadercolor}{rgb}{0.709,0.784,0.454}
				\definecolor{tablerowAltcolor}{rgb}{.866,.905,.737}
				\definecolor{tablerowAlt2color}{rgb}{.968,.976,.933}
				\definecolor{verylightgray}{rgb}{0.98,0.98,0.98}
				
				\newenvironment{fshaded}{
				\def\FrameCommand{\fcolorbox{framecolor}{shadecolor}}
				\MakeFramed {\FrameRestore}}
				{\endMakeFramed}
				
				\newenvironment{fexample}[1][]{\definecolor{shadecolor}{rgb}{.913,.913,.913}
				\definecolor{framecolor}{rgb}{.156,.333,.443}
				\begin{fshaded}}{\end{fshaded}} 
				
				\newenvironment{fdefinition}{\definecolor{shadecolor}{rgb}{.913,.913,.913}
				\definecolor{framecolor}{rgb}{.709,.784,.454}
				\begin{fshaded}}{\end{fshaded}}
				
				\newenvironment{fexercise}{\definecolor{shadecolor}{rgb}{.913,.913,.913}
				\definecolor{framecolor}{rgb}{.49,.639,0}
				\begin{fshaded}}{\end{fshaded}}
				
				\newenvironment{fhint}{\definecolor{shadecolor}{rgb}{.913,.913,.913}
				\definecolor{framecolor}{rgb}{.941,.674,.196}
				\begin{fshaded}}{\end{fshaded}}	
				
				\newcommand{\PreserveBackslash}[1]{
				\let\temp=\\#1\let\\=\temp
				}
				\let\PBS=\PreserveBackslash
				\newcolumntype{A}{>{\PBS\raggedright\small\hspace{0pt}}X}
				\newcolumntype{L}[1]{>{\PBS\raggedright\small\hspace{0pt}}p{#1}}
				\newcolumntype{R}[1]{>{\PBS\raggedleft\small\hspace{0pt}}p{#1}}
				\newcolumntype{C}[1]{>{\PBS\centering\small\hspace{0pt}}p{#1}}
				
				\makeindex
				
				\title{TD S\'equentiel - Rappels de base}	
			\date{}
			\author{\scriptsize{}}
			\definecolor{light-gray}{gray}{0.8}
			\renewcommand{\headrulewidth}{0pt}
			\fancyhead[L]{
				\footnotesize\textsc{Haute \'Ecole de Bruxelles}\\
	    			\footnotesize\textsc{\'Ecole Sup\'erieure d'Informatique}
			}
			\fancyhead[R]{
				\footnotesize{Bachelor en Informatique}\\
				\footnotesize{Laboratoires Java} - 
			\footnotesize{1\`ere ann\'ee}}
				\fancyfoot[L]{ }
				\fancyfoot[C]{}
				\fancyfoot[R]{\scriptsize{\textcolor{gray}{version 2014-2015 (\today)}}}
				\pagestyle{plain}
				\reversemarginpar
				\usepackage{rotating}						
				\begin{document}
					\begin{textblock}{9}(2,3.2)
						\includegraphics[width=2cm]{../../../_templates/java/icons/logo-esi}
					\end{textblock}
				
				
				
				
				%\maketitle
				\cadretitre{TD1}{TD S\'equentiel - Rappels de base}
				\thispagestyle{fancy}
        \marginpar{\begin{sideways}
            \begin{minipage}[t]{1cm}
            \begin{tiny}
            \includegraphics[width=1\linewidth,height=1\textheight,keepaspectratio=true]{../../../_templates/java/icons/cc-gris.jpg}
			\end{tiny}
			\end{minipage}
            \begin{minipage}[b]{19cm}
            \begin{tiny}
            \textcolor{gray}{Distribué sous licence Creative Commons Paternité - Partage à l'Identique 2.0 Belgique 
            (\texttt{http://creativecommons.org/licenses/by-sa/2.0/be/})
			\vspace{-1em}
			\\Les autorisations au-delà du champ de cette licence peuvent être obtenues à 
			\texttt{http://www.heb.be/esi}
			- \texttt{mcodutti@heb.be}
			}\end{tiny}
			\end{minipage}
        \end{sideways}}
            \begin{abstract}
			Ce TD a pour but de 
			
            \par
        
            \par
        \end{abstract}
				\vspace{-2em}\tableofcontents
				\pagestyle{plain}
            \clearpage
            \fancyhead[L,C,R]{}
            \fancyfoot[L,C]{}
            \fancyfoot[R]{ \scriptsize{\textcolor{gray}{
				TDSeq - page \thepage}}}
				\thispagestyle{fancy}
				\pagestyle{fancy}
	   
            \section{Algorithmes s\'equentiels}Revoyons ici les bases du pseudo-code et leur traduction en Java.\subsection{Le pseudo-code}
		    Le pseudo-code ou Langage de Description des Algorithmes (LDA en abr\'eg\'e) est un langage
        formel et symbolique utilisant :
        
					\begin{itemize}
				
			\item 
            des noms symboliques destin\'es \`a repr\'esenter les objets sur lesquels s'effectuent des
            actions ;
          
			\item 
            des op\'erateurs symboliques ou des mots-cl\'es traduisant les op\'erations primitives
            ex\'ecutables par un ex\'ecutant donn\'e ;
          
			\item 
            des structures de contr\^ole types.
          
					\end{itemize}
				
            \par
        \section{Variables et types}
				Nous savons que les op\'erations que l'ordinateur devra ex\'ecuter portent sur des \'el\'ements qui
        sont les \textbf{donn\'ees} du probl\`eme.
      
            \par
        
        Lorsqu'on attribue un \textbf{nom} 
        et un \textbf{type} \`a ces donn\'ees, on
        parle alors de \textbf{variables}. 
      
            \par
        
        Dans un algorithme, une variable conserve toujours son nom et
        son type, mais peut changer de \textbf{valeur}.
         
					\begin{itemize}
				
			\item 
            Le \textbf{nom} d'une variable permet de la caract\'eriser et de la reconnaitre ;
          
			\item 
            le \textbf{type} d'une variable d\'ecrit la nature de son contenu.
          
					\end{itemize}
				
            \par
        \subsection{ Les types autoris\'es en algo}
          Dans un premier temps, les seuls types utilis\'es sont :
          
					\begin{itemize}
				
			\item \verb@entier@ pour les nombres entiers ;
			\item \verb@réel@ pour les nombres r\'eels ;
			\item \verb@caractère@ pour les diff\'erentes lettres et caract\`eres
              (par exemple ceux qui apparaissent sur un clavier : 'a', '1', '\#', etc.)
			\item \verb@chaine@ pour les variables contenant un ou plusieurs caract\`ere(s) ou aucun (la chaine vide)
              (par exemple : "Bonjour", "Bonjour le monde", "a", "", etc.)
			\item \verb@booléen@ les variables de ce type 
            ne peuvent valoir que \verb@vrai@ ou \verb@faux@
					\end{itemize}
				
            \par
        
			
		\subparagraph{Les commandes de base} 
		
                \textcolor{white}{.} \par
             
								Quel(s) type(s) de donn\'ees utiliseriez-vous pour repr\'esenter :
							
					\begin{itemize}
				
			\item une date du calendrier ?  \textcolor{gray}{\underline{\hspace*{10em}}} 
			\item un moment dans la journ\'ee ?  \textcolor{gray}{\underline{\hspace*{10em}}} 
			\item le prix d'un produit en grande surface ? \textcolor{gray}{\underline{\hspace*{5em}}} 
			\item votre nom ?  \textcolor{gray}{\underline{\hspace*{10em}}} 
			\item vos initiales ? \textcolor{gray}{\underline{\hspace*{10em}}} 
			\item votre adresse ?  \textcolor{gray}{\underline{\hspace*{10em}}} 
					\end{itemize}
				\subsection{ Les types \'equivalents en java}
          Dans un premier temps, les seuls types utilis\'es sont :
          
					\begin{itemize}
				
			\item \verb@entier@ pour les nombres entiers ;
			\item \verb@réel@ pour les nombres r\'eels ;
			\item \verb@caractère@ pour les diff\'erentes lettres et caract\`eres
              (par exemple ceux qui apparaissent sur un clavier : 'a', '1', '\#', etc.)
			\item \verb@chaine@ pour les variables contenant un ou plusieurs caract\`ere(s) ou aucun (la chaine vide)
              (par exemple : "Bonjour", "Bonjour le monde", "a", "", etc.)
			\item \verb@booléen@ les variables de ce type 
            ne peuvent valoir que \verb@vrai@ ou \verb@faux@
					\end{itemize}
				
            \par
        
			
		\subparagraph{Les commandes de base} 
		
                \textcolor{white}{.} \par
             
								Quel(s) type(s) de donn\'ees utiliseriez-vous pour repr\'esenter :
							
					\begin{itemize}
				
			\item une date du calendrier ?  \textcolor{gray}{\underline{\hspace*{10em}}} 
			\item un moment dans la journ\'ee ?  \textcolor{gray}{\underline{\hspace*{10em}}} 
			\item le prix d'un produit en grande surface ? \textcolor{gray}{\underline{\hspace*{5em}}} 
			\item votre nom ?  \textcolor{gray}{\underline{\hspace*{10em}}} 
			\item vos initiales ? \textcolor{gray}{\underline{\hspace*{10em}}} 
			\item votre adresse ?  \textcolor{gray}{\underline{\hspace*{10em}}} 
					\end{itemize}
				
				\end{document}
			